\documentclass[12pt,a4paper]{article}

% ─── Pacchetti ──────────────────────────────────────────────────────────────
\usepackage[T1]{fontenc}
\usepackage[utf8]{inputenc}
\usepackage[italian]{babel}
\usepackage[margin=2.5cm]{geometry}
\usepackage{hyperref}
\usepackage{amsmath,amssymb}
\usepackage{listings}
\usepackage{xcolor}
\usepackage{graphicx}
\usepackage{enumitem}
\usepackage{booktabs}
\usepackage{microtype}
\usepackage{fancyhdr}
\usepackage{titlesec}
\usepackage{abstract}
\usepackage{setspace}

% ─── Stile pagina ────────────────────────────────────────────────────────────
\pagestyle{fancy}
\fancyhf{}
\rhead{\textit{Padding Oracle SOC Lab}}
\lhead{\textit{Specifiche Accademiche}}
\cfoot{\thepage}

% ─── Titoli sezione ─────────────────────────────────────────────────────────
\titleformat{\section}{\large\bfseries}{\thesection}{1em}{}[\titlerule]
\titleformat{\subsection}{\normalsize\bfseries}{\thesubsection}{1em}{}
\titleformat{\subsubsection}{\normalsize\itshape}{\thesubsubsection}{1em}{}

% ─── Stile codice ────────────────────────────────────────────────────────────
\definecolor{codebg}{RGB}{248,248,248}
\definecolor{codeframe}{RGB}{200,200,200}
\definecolor{keyword}{RGB}{0,0,180}
\definecolor{comment}{RGB}{80,120,80}
\definecolor{string}{RGB}{160,32,32}

\lstset{
  backgroundcolor=\color{codebg},
  frame=single,
  rulecolor=\color{codeframe},
  basicstyle=\ttfamily\footnotesize,
  keywordstyle=\color{keyword}\bfseries,
  commentstyle=\color{comment}\itshape,
  stringstyle=\color{string},
  numbers=left,
  numberstyle=\tiny\color{gray},
  numbersep=6pt,
  breaklines=true,
  showstringspaces=false,
  captionpos=b,
}

% ─── Frontespizio ────────────────────────────────────────────────────────────
\title{
  \vspace{-1cm}
  {\Large Corso di Sicurezza Informatica}\\[0.4cm]
  {\rule{\linewidth}{0.5pt}}\\[0.4cm]
  {\LARGE\bfseries Padding Oracle SOC Lab}\\[0.2cm]
  {\large Progettazione e Implementazione di un Ambiente\\
   Didattico per l'Analisi di Attacchi Crittografici\\
   e il Rilevamento Automatico tramite SIEM}\\[0.4cm]
  {\rule{\linewidth}{0.5pt}}
}
\author{
  Studente: \textit{[Nome Cognome]}\\
  Matricola: \textit{[000000]}\\[0.2cm]
  Relatore: \textit{[Prof. Nome Cognome]}\\
  Dipartimento di Informatica
}
\date{Anno Accademico 2025--2026 \\ \vspace{0.2cm} \small Commit \texttt{80982df} — branch \texttt{master}}

\begin{document}

\maketitle
\thispagestyle{empty}
\clearpage

% ─── Abstract ────────────────────────────────────────────────────────────────
\begin{abstract}
\noindent
Il presente documento descrive la progettazione, la motivazione teorica e l'architettura del
\emph{Padding Oracle SOC Lab}, un ambiente containerizzato sviluppato per lo studio accademico
degli attacchi crittografici basati su \emph{Padding Oracle} in modalità \mbox{AES-CBC} e per
la sperimentazione di un sistema di rilevamento automatico delle intrusioni di tipo SIEM
(\emph{Security Information and Event Management}).
Il laboratorio consente di riprodurre in modo controllato tre scenari di esposizione — vulnerabilità
esplicita, side-channel temporale e mitigazione completa — e di osservare in tempo reale come un
motore di correlazione statistica sia in grado di distinguere traffico malevolo da traffico legittimo,
calcolare metriche formali di detection e attivare risposte automatiche tramite integrazione SOAR.
\end{abstract}

\tableofcontents
\clearpage

% ===========================================================================
\section{Introduzione e Motivazioni}
% ===========================================================================

\subsection{Contesto}

La crittografia simmetrica a blocchi costituisce uno dei pilastri della sicurezza applicativa moderna.
Algoritmi come l'\emph{Advanced Encryption Standard} (AES), standardizzato dal NIST nel 2001
\cite{nist-aes}, sono impiegati in milioni di sistemi per proteggere sessioni, token e dati a riposo.
Tuttavia, la correttezza di un algoritmo crittograficamente sicuro non implica necessariamente la
sicurezza del \emph{protocollo} che lo incapsula: errori di design a livello applicativo possono
esporre canali di informazione indiretti, noti in letteratura come \emph{oracle}.

Il \textbf{Padding Oracle Attack} — formalizzato da Vaudenay nel 2002 \cite{vaudenay2002} e
successivamente esteso da Rizzo e Duong con gli attacchi \emph{BEAST} (2011) e
\emph{POODLE} (2014) \cite{poodle2014} — rappresenta uno degli esempi più pedagogicamente
rilevanti di come una primitiva sicura possa essere resa vulnerabile da un contesto applicativo
che divulga informazioni sullo stato interno del cifrario.

\subsection{Obiettivi del Laboratorio}

Il laboratorio si propone di raggiungere i seguenti obiettivi accademici:

\begin{enumerate}[label=\textbf{O\arabic*.}]
  \item \textbf{Comprensione teorica} della modalità CBC e del meccanismo di padding PKCS\#7,
        con particolare attenzione alle condizioni che rendono l'oracolo sfruttabile.
  \item \textbf{Implementazione pratica} di un exploit funzionante, capace di decifrare token
        AES-CBC senza conoscenza della chiave segreta.
  \item \textbf{Studio dei canali laterali} (\emph{side-channel}): come differenze di latenza
        nell'elaborazione crittografica possono essere sfruttate anche in assenza di messaggi
        di errore espliciti.
  \item \textbf{Progettazione e validazione} di un motore di rilevamento SIEM basato su
        metriche statistiche, capace di distinguere traffico di attacco da traffico benigno
        con alta precisione e basso tasso di falsi positivi.
  \item \textbf{Sperimentazione delle mitigazioni}: dalla correzione architetturale
        (\emph{Encrypt-then-MAC}) alla difesa perimetrale inline (WAF con sliding-window).
\end{enumerate}

% ===========================================================================
\section{Fondamenti Teorici: AES in Modalità CBC}
% ===========================================================================

\subsection{Block Cipher e Modalità Operative}

Un \emph{block cipher} è una funzione biettiva $E_k : \{0,1\}^n \to \{0,1\}^n$ parametrizzata
da una chiave segreta $k$. AES opera su blocchi di 128 bit ($n = 128$) con chiavi di 128, 192
o 256 bit. Poiché i messaggi reali hanno lunghezza arbitraria, si rende necessaria una
\emph{modalità operativa} che gestisca la cifratura di sequenze di blocchi.

\subsection{Cipher Block Chaining (CBC)}

La modalità \emph{Cipher Block Chaining} (CBC), definita in \cite{nist-sp800-38a}, introduce
una dipendenza sequenziale tra i blocchi. Dati i blocchi di plaintext
$P_1, P_2, \ldots, P_m$ e un vettore di inizializzazione casuale $IV$, la cifratura è definita da:

\begin{equation}
  C_i = E_k\!\left(P_i \oplus C_{i-1}\right), \quad C_0 = IV
\end{equation}

La corrispondente decifratura è:

\begin{equation}
  P_i = D_k(C_i) \oplus C_{i-1}
\end{equation}

dove $D_k$ è la funzione di decifratura AES con chiave $k$. Si osservi che la decifratura
di ogni blocco è indipendente dalla chiave una volta noto $D_k(C_i)$, termine che dipende
esclusivamente dal testo cifrato e dalla chiave — ma non dalla conoscenza esplicita di $k$
da parte dell'attaccante.

\subsection{Padding PKCS\#7}

Poiché i messaggi reali raramente hanno lunghezza multipla della dimensione del blocco
($\ell = 16$ byte per AES), è necessario un meccanismo di \emph{padding} per portare
il plaintext alla dimensione corretta. Lo standard PKCS\#7 \cite{pkcs7} stabilisce che,
se mancano $p$ byte all'allineamento al blocco, si aggiungono esattamente $p$ byte ciascuno
con valore $p$:

\begin{equation}
  \text{pad}(M) = M \,\|\, \underbrace{p \,\|\, p \,\|\, \cdots \,\|\, p}_{p \text{ volte}},
  \quad p = \ell - (|M| \bmod \ell)
\end{equation}

con $p \in \{1, 2, \ldots, 16\}$. La verifica del padding in fase di decifratura consiste
nel leggere il valore dell'ultimo byte $v$ e verificare che gli ultimi $v$ byte siano tutti
uguali a $v$. Un'implementazione errata che restituisce risposte \emph{differenziate} in
caso di padding valido o invalido costituisce il requisito fondamentale per un Padding Oracle.

% ===========================================================================
\section{Il Padding Oracle Attack}
% ===========================================================================

\subsection{Definizione Formale di Oracle}

Nella teoria della complessità computazionale, un \emph{oracolo} è un dispositivo che risponde
a query in tempo costante. Nel contesto crittografico, un \emph{Padding Oracle} è qualsiasi
entità — tipicamente un server — che, dato un ciphertext $C$, è in grado di comunicare
(esplicitamente o implicitamente) se il plaintext decifrato soddisfa o meno il formato di
padding atteso.

\begin{definition}[Padding Oracle]
Sia $D_k : \{0,1\}^* \to \{0,1\}^*$ una funzione di decifratura con padding. Un
\emph{Padding Oracle} $\mathcal{O}$ è una funzione booleana:
\[
  \mathcal{O}(C) = \begin{cases} 1 & \text{se } \mathrm{unpad}(D_k(C)) \neq \bot \\ 0 & \text{altrimenti} \end{cases}
\]
\end{definition}

L'esistenza di $\mathcal{O}$ trasforma un cifrario IND-CPA sicuro in un sistema completamente
insicuro sotto attacchi adattativi con testo scelto (IND-CCA), come dimostrato formalmente
da Vaudenay \cite{vaudenay2002}.

\subsection{Meccanismo dell'Attacco: Decifratura Byte per Byte}

L'attacco sfrutta la struttura CBC e la linearità dell'operazione XOR per recuperare il
plaintext byte per byte, senza mai conoscere la chiave $k$.

\subsubsection{Analisi di un Singolo Byte}

Si consideri l'attacco sull'ultimo byte del blocco $P_i$. L'attaccante conosce $C_{i-1}$
e $C_i$, e vuole determinare $P_i[15]$ (l'ultimo byte, con indicizzazione da 0).

Per la proprietà della decifratura CBC:
\[
  P_i[15] = D_k(C_i)[15] \oplus C_{i-1}[15]
\]

L'attaccante costruisce un blocco modificato $C'_{i-1}$ variando solo l'ultimo byte:
\[
  C'_{i-1}[15] = g \quad (\text{con } g \in \{0, 1, \ldots, 255\})
\]

e interroga l'oracolo con il ciphertext $(C'_{i-1}, C_i)$. Quando $\mathcal{O}(C'_{i-1}, C_i) = 1$,
il padding risultante è \texttt{0x01}, il che implica:
\[
  D_k(C_i)[15] \oplus g^* = \texttt{0x01}
  \implies D_k(C_i)[15] = g^* \oplus \texttt{0x01}
  \implies P_i[15] = D_k(C_i)[15] \oplus C_{i-1}[15]
\]

dove $g^*$ è il valore che ha provocato una risposta positiva dell'oracolo.

\subsubsection{Generalizzazione a Tutti i Byte}

Per recuperare il byte in posizione $j$ (con $j = 15, 14, \ldots, 0$), si imposta il padding
target a $p = 16 - j$ e si aggiustano i byte già recuperati in modo che soddisfino il padding
desiderato. Il processo richiede al più $256 \cdot 16 = 4096$ query per blocco
(nella pratica molte meno grazie alla distribuzione uniforme dei byte di padding valido).

\subsubsection{Complessità dell'Attacco}

Per un messaggio di $m$ blocchi da 16 byte ($m \cdot 16 \cdot 8 = 128m$ bit totali):
\begin{itemize}
  \item \textbf{Query nel caso peggiore:} $255 \cdot 16 \cdot m = 4080m$
  \item \textbf{Query attese:} $128m$ (un tentativo per byte in media con distribuzione uniforme)
  \item \textbf{Complessità:} $O(n)$ in $n = |$plaintext$|$ — efficienza polinomiale, non esponenziale
\end{itemize}

Il contrasto con la sicurezza IND-CCA è netto: un avversario con accesso all'oracolo può
decifrare un qualsiasi messaggio in tempo lineare nella lunghezza del ciphertext.

\subsection{Variante Side-Channel: Timing Oracle}

In assenza di risposte di errore esplicite, un attaccante sofisticato può sfruttare
differenze di \emph{latenza} nel tempo di elaborazione del server. Questo scenario, noto
come \emph{Timing Oracle}, si verifica quando:

\begin{itemize}
  \item La verifica del padding fallisce \emph{prima} rispetto a quando ha successo
        (early exit nel codice di validazione), oppure
  \item Un'implementazione non \emph{constant-time} introduce rami condizionali visibili
        come differenze di latenza statisticamente significative.
\end{itemize}

La rilevazione di un Timing Oracle richiede analisi statistica delle latenze:
si confrontano le distribuzioni dei tempi di risposta per richieste con padding valido
versus invalido. Una distribuzione \emph{bimodale} (due picchi distinti) è la firma
caratteristica del canale laterale temporale.

\subsection{Attacchi Reali Notevoli}

\begin{description}
  \item[POODLE (2014)] \hfill \cite{poodle2014}\\
    Vulnerabilità nel protocollo SSLv3 che sfrutta un padding oracle per decifrare
    traffico HTTPS. Resa pubblica da Möller, Duong e Kotowicz di Google Security Team.
  \item[Lucky Thirteen (2013)] \hfill \cite{lucky13}\\
    Attacco timing-based contro TLS 1.0--1.2 che sfrutta differenze di latenza
    nel codice MAC di OpenSSL, con misurazioni nell'ordine dei nanosecondi.
  \item[BEAST (2011)] \hfill \cite{beast2011}\\
    Attacco chosen-plaintext in CBC che sfrutta la prevedibilità dell'IV in
    SSL/TLS 1.0 per decifrare sessioni HTTPS.
\end{description}

% ===========================================================================
\section{Descrizione del Progetto}
% ===========================================================================

\subsection{Visione d'Insieme e Finalità}

Il \emph{Padding Oracle SOC Lab} è un sistema didattico progettato per colmare il divario
tra la teoria formale degli attacchi crittografici e la loro rilevazione operativa in un
contesto di sicurezza realistico. A differenza dei classici laboratori che presentano
solo l'exploit o solo la difesa, questo progetto integra entrambe le prospettive in un
singolo ambiente coerente, consentendo di osservare l'intero ciclo di vita di un attacco:
dalla sua esecuzione, alla sua rilevazione tramite telemetria, fino alla risposta automatica.

\subsection{Architettura del Sistema}

Il sistema adotta un'architettura a micro-servizi containerizzati che comunicano tramite
uno strato di \emph{event-driven file streaming}: ogni componente produce log strutturati
in formato JSONL (\emph{JSON Lines}) in una directory condivisa, che il motore SIEM
consuma in modo incrementale tramite lettura con offset di byte. Questa scelta elimina la
dipendenza da un message broker esterno, mantenendo il sistema autonomo e reproducibile.

\subsubsection{Componente Victim}

Il \emph{victim} è un'applicazione web Flask che espone endpoint crittografici. Il token
ha la struttura:

\[
  \text{Token} = IV \,\|\, E_k\!\bigl(\text{pad}(\text{plaintext} \,\|\, \text{HMAC}_{k'}(\text{plaintext}))\bigr)
\]

dove $k$ è la chiave AES-128-CBC e $k'$ la chiave HMAC-SHA256, entrambe parametrizzabili
tramite variabili d'ambiente. Il victim è configurabile in tre modalità operative:

\begin{enumerate}
  \item \textbf{Vulnerabile (\texttt{vuln})}: In caso di padding invalido, il server risponde
        con HTTP 500 e un corpo JSON contenente il campo \texttt{"result": "padding\_error"}.
        In caso di errore di integrità HMAC, risponde con HTTP 403 e
        \texttt{"result": "integrity\_error"}. Questa distinzione costituisce un Padding Oracle
        classico e diretto.

  \item \textbf{Parzialmente mitigato (\texttt{partial})}: Il server applica un ritardo
        artificiale differenziato: 30 ms per errori di padding (simulando elaborazione
        crittografica più lenta), 5 ms per errori di integrità. Il corpo della risposta
        è unificato (\texttt{"result": "request\_denied"}) ma il canale laterale temporale
        rimane sfruttabile.

  \item \textbf{Corretto (\texttt{fixed})}: Implementazione \emph{Encrypt-then-MAC} con
        verifica dell'integrità prima della decifratura, ritardo costante di 15 ms
        indipendente dal tipo di errore, e risposta unificata. Elimina sia il Padding Oracle
        esplicito sia il Timing Oracle.
\end{enumerate}

\subsubsection{Componente Attacker}

Lo script di attacco implementa l'algoritmo Padding Oracle descritto nella
Sezione~\ref{sec:meccanismo} in Python. Opera iterando sui blocchi del ciphertext da destra
a sinistra, interrogando sistematicamente l'endpoint \texttt{/decrypt} del victim con
variazioni del blocco $C_{i-1}$. Ogni query viene annotata con metadati (IP sorgente,
timestamp, lunghezza del ciphertext, byte target) e trasmessa al SOC tramite log JSONL.

\subsubsection{Componente Benign Client}

Il client benigno simula traffico legittimo multi-IP con pattern di accesso variabile
(burst casuale, rate stazionario). Il suo scopo è fornire una \emph{baseline} realistica
al motore di detection, consentendo di misurare il tasso di falsi positivi quando le
regole di rilevamento operano in presenza di traffico normale.

\subsubsection{Componente SOC Collector e Dashboard}

Il motore SOC è l'elemento centrale del sistema di difesa. Aggrega gli eventi JSONL
da tutti i componenti, li persiste in un database SQLite con modalità WAL
(\emph{Write-Ahead Logging}) per ottimizzare le letture concorrenti, e applica le regole
di detection descritte nella Sezione~\ref{sec:siem}. La dashboard web offre visualizzazione
in tempo reale della topologia di rete, degli alert attivi, delle metriche KPI e delle
configurazioni WAF, utilizzando la libreria \texttt{vis-network.js} per il grafo.

\subsection{Infrastruttura e Deploy}

Il sistema è orchestrato tramite Docker Compose su una rete bridge isolata. I servizi
espongono le seguenti porte sull'host:

\begin{center}
\begin{tabular}{lll}
\toprule
\textbf{Servizio} & \textbf{Porta Host} & \textbf{Funzione} \\
\midrule
\texttt{victim}    & 18080 & API crittografiche (target dell'attacco) \\
\texttt{soc}       & 18090 & API REST SIEM (\texttt{/alerts}, \texttt{/metrics}) \\
\texttt{soc-ui}    & 18091 & Dashboard web interattiva \\
\bottomrule
\end{tabular}
\end{center}

% ===========================================================================
\section{SIEM: Architettura e Regole di Rilevamento}
\label{sec:siem}
% ===========================================================================

\subsection{Fondamenti di SIEM e SOC}

Un sistema SIEM (\emph{Security Information and Event Management}) ha il compito di
raccogliere, normalizzare, correlare e analizzare log di sicurezza provenienti da sorgenti
eterogenee, al fine di rilevare anomalie e generare alert azionabili. La sua efficacia
è misurata attraverso metriche formali standard:

\begin{description}
  \item[\textbf{TPR} (True Positive Rate):] Frazione di attacchi reali correttamente rilevati.
        $\text{TPR} = \frac{TP}{TP + FN}$
  \item[\textbf{FPR} (False Positive Rate):] Frazione di eventi benigni erroneamente classificati
        come malevoli. $\text{FPR} = \frac{FP}{FP + TN}$
  \item[\textbf{MTTD} (Mean Time to Detect):] Tempo medio tra l'inizio di un attacco e la
        generazione del primo alert. Misura la reattività del sistema.
\end{description}

\subsection{Pipeline di Ingestion e Correlazione}

La pipeline del SOC si articola in tre fasi:

\begin{enumerate}
  \item \textbf{Ingestion incrementale:} Il modulo \texttt{collector.py} legge i file JSONL
        con un meccanismo a offset di byte (\emph{tail-like}), garantendo che nessun evento
        venga processato due volte. I nuovi eventi vengono inseriti nel database SQLite tramite
        una transazione atomica con lock dedicato.

  \item \textbf{Normalizzazione:} Ogni evento viene normalizzato verso uno schema comune
        che include: \texttt{event\_id}, \texttt{timestamp} (ISO 8601), \texttt{src\_ip},
        \texttt{endpoint}, \texttt{status\_code}, \texttt{latency\_ms},
        \texttt{crypto\_time\_ns}, \texttt{error\_type} e \texttt{mode}.

  \item \textbf{Correlazione statistica:} Il motore applica regole basate su aggregazioni
        per IP sorgente su finestre temporali configurabili (\emph{sliding window}).
\end{enumerate}

\subsection{Regole di Rilevamento}

Il sistema implementa tre regole principali, ciascuna associata a una tecnica del framework
MITRE ATT\&CK:

\subsubsection{Regola 1 — Error Flooding (T1110.001)}

Questa regola identifica il Padding Oracle classico nella modalità \texttt{vuln}, dove
le risposte differenziate sono esplicite. Il trigger si attiva quando:

\[
  N_{IP} \geq \tau_{min} \quad \land \quad
  \frac{\#\{\text{errori}\}}{N_{IP}} \geq \tau_{fail} \quad \land \quad
  \#\{\text{padding\_error}\} > 0
\]

con soglie di default $\tau_{min} = 25$ eventi e $\tau_{fail} = 0.85$.
La presenza di almeno un evento \texttt{padding\_error} (HTTP 500) è il discriminatore
che distingue questa regola da un generico errore applicativo.

\subsubsection{Regola 2 — Timing Side-Channel (T1595.002)}

Questa regola rileva il Timing Oracle nella modalità \texttt{partial} attraverso
l'analisi della distribuzione delle latenze. Il detector calcola il
\emph{coefficiente di bimodalità di Sarle} (BC), una misura statistica che quantifica
quanto una distribuzione si discosta dalla normalità verso due picchi distinti:

\[
  BC = \frac{\hat{\gamma}_1^2 + 1}{\hat{\kappa} + \frac{3(n-1)^2}{(n-2)(n-3)}}
\]

dove $\hat{\gamma}_1$ è il coefficiente di asimmetria campionario e $\hat{\kappa}$
la curtosi di Pearson. Per una distribuzione normale, $BC \approx 0.333$; una
distribuzione bimodale produce $BC > 0.555$, che costituisce la soglia di alert.

Il trigger si attiva quando almeno una delle seguenti condizioni è soddisfatta
(con $n \geq 20$ campioni):
\begin{itemize}
  \item $\sigma_{latency} \geq 6.0\,\text{ms}$
  \item $P_{95} - P_{50} \geq 12.0\,\text{ms}$
  \item $BC \geq 0.555$
\end{itemize}

\subsubsection{Regola 3 — Block Probing Signature (T1040)}

Questa regola riconosce il pattern strutturale dell'exploit: l'attaccante invia
sequenze di richieste con ciphertext della \emph{stessa lunghezza} (multiplo di 16,
cioè allineata al blocco AES) ma con contenuto sistematicamente variato. Il trigger è:

\[
  \#\{\text{errori}\} \geq \tau_{consec} \quad \land \quad
  \text{fail\_rate} \geq 0.70 \quad \land \quad
  |\mathcal{L}| = 1 \quad \land \quad
  \ell_0 \equiv 0 \pmod{16}
\]

dove $\mathcal{L}$ è l'insieme delle lunghezze di ciphertext osservate, $\ell_0$
l'unica lunghezza presente, e $\tau_{consec} = 15$ il numero minimo di errori
consecutivi.

\subsection{WAF Inline e Risposta Automatica (SOAR)}

Il sistema integra un \emph{Web Application Firewall} inline implementato direttamente
nel victim, che valuta ogni richiesta in ingresso rispetto a una policy di sliding-window
per IP sorgente. Due regole principali sono pre-configurate:

\begin{description}
  \item[WAF-01 (Error Flooding \& High Fail-Rate Ban):] Blocca l'IP quando il fail-rate
        supera l'80\% con almeno 15 richieste in finestra di 60 secondi. Ban TTL: 120 s.
  \item[WAF-02 (Consecutive CBC Byte-Probing):] Blocca l'IP dopo 8 errori crittografici
        consecutivi in 60 secondi. Ban TTL: 180 s.
\end{description}

Il componente SOAR (\emph{Security Orchestration, Automation and Response}) consente
al SIEM di richiedere automaticamente il blocco dell'IP malevolo sul WAF del victim
tramite una chiamata REST asincrona, chiudendo il ciclo di risposta senza intervento
umano.

\subsection{Generazione di Regole Sigma}

Il sistema è in grado di esportare automaticamente le regole di detection nel formato
Sigma \cite{sigma}, uno standard aperto per la scrittura di regole SIEM indipendenti
dalla piattaforma. Le regole generate includono i metadati MITRE ATT\&CK, le condizioni
di trigger e i campi di evidenza, e sono pronte per essere importate in piattaforme
come Splunk, Elasticsearch o Microsoft Sentinel.

% ===========================================================================
\section{Sezione Pratica: Esecuzione e Verifica}
% ===========================================================================

Questa sezione, unica parte operativa del documento, riporta i comandi necessari per
avviare il laboratorio, riprodurre l'attacco e verificare il funzionamento del SIEM.
Per i dettagli implementativi si rimanda al codice sorgente.

\begin{lstlisting}[language=bash, caption={Avvio completo del laboratorio e riproduzione dell'attacco}]
# 1. Avvio dello stack completo (victim in modalita' vuln, soc, soc-ui)
./control/scenario.sh up

# 2. Verifica che tutti i servizi siano operativi
curl http://localhost:18080/health   # victim
curl http://localhost:18090/health   # soc collector
# Aprire http://localhost:18091 nel browser per la dashboard

# 3. Avvio del traffico benigno (necessario per costruire la baseline SIEM)
./control/scenario.sh start-benign

# 4. Esecuzione dell'attacco Padding Oracle contro il victim vulnerabile
./control/scenario.sh run-attack
# oppure direttamente:
python attacker/attack.py --target http://localhost:18080

# 5. Abilitare le regole SIEM e osservare gli alert generati
curl -X POST http://localhost:18090/alerts/rules \
  -H "Content-Type: application/json" \
  -d '{"enabled": true}'
curl http://localhost:18090/alerts   # visualizza alert attivi

# 6. Passare alla modalita' timing side-channel e ripetere
./control/scenario.sh up-partial
./control/scenario.sh run-attack

# 7. Verificare che la modalita' fixed blocchi l'attacco
./control/scenario.sh up-fixed
./control/scenario.sh run-attack    # deve fallire il recupero del plaintext

# 8. Eseguire la suite di test completa
pytest tests/ -v
\end{lstlisting}

% ===========================================================================
\section{Conclusioni}
% ===========================================================================

Il \emph{Padding Oracle SOC Lab} costituisce un caso di studio completo che abbraccia
tre dimensioni complementari della sicurezza informatica: la \emph{crittografia applicata}
(AES-CBC, PKCS\#7, HMAC), la \emph{sicurezza offensiva} (exploit padding oracle e
timing side-channel) e la \emph{sicurezza difensiva} (detection statistica, WAF, SOAR).

La scelta di implementare tre modalità operative del victim non è casuale: riflette
la progressione naturale delle mitigazioni osservata storicamente nell'industria,
da sistemi completamente vulnerabili a soluzioni parziali (che eliminano il canale
esplicito ma espongono quello temporale) fino alla correzione architetturale corretta
(\emph{Encrypt-then-MAC} con costant-time comparison).

Il motore SIEM dimostra che è possibile rilevare entrambe le varianti dell'attacco
— esplicita e temporale — attraverso metriche statisticamente fondate (bimodality
coefficient, standard deviation, fail rate), senza richiedere firme statiche o
pattern di payload, e quindi resistenti a variazioni superficiali dell'exploit.

% ===========================================================================
% Riferimenti
% ===========================================================================
\begin{thebibliography}{99}

\bibitem{vaudenay2002}
  S.~Vaudenay,
  ``Security Flaws Induced by CBC Padding — Applications to SSL, IPSEC, WTLS...'',
  in \textit{Advances in Cryptology — EUROCRYPT 2002},
  Lecture Notes in Computer Science, vol.~2332, Springer, 2002.

\bibitem{nist-aes}
  National Institute of Standards and Technology,
  \textit{FIPS PUB 197: Advanced Encryption Standard (AES)},
  Federal Information Processing Standards Publication, 2001.

\bibitem{nist-sp800-38a}
  M.~Dworkin,
  \textit{NIST Special Publication 800-38A: Recommendation for Block Cipher Modes of Operation},
  National Institute of Standards and Technology, 2001.

\bibitem{pkcs7}
  B.~Kaliski,
  \textit{RFC 2315: PKCS \#7: Cryptographic Message Syntax Version 1.5},
  Internet Engineering Task Force, 1998.

\bibitem{poodle2014}
  B.~Möller, T.~Duong, K.~Kotowicz,
  ``This POODLE Bites: Exploiting The SSL 3.0 Fallback'',
  Google Security Advisory, settembre 2014.

\bibitem{lucky13}
  N.~J.~Al~Fardan, K.~G.~Paterson,
  ``Lucky Thirteen: Breaking the TLS and DTLS Record Protocols'',
  in \textit{IEEE Symposium on Security and Privacy}, 2013.

\bibitem{beast2011}
  J.~Rizzo, T.~Duong,
  ``Here Come the XOR Ninjas'',
  ekoparty Security Conference, 2011.

\bibitem{sigma}
  T.~Patzke et al.,
  ``Sigma: Generic Signature Format for SIEM Systems'',
  \url{https://github.com/SigmaHQ/sigma}, 2017.

\end{thebibliography}

\vfill
\begin{center}
\small\textit{Generato da \texttt{/project-brief} — commit \texttt{80982df} — branch \texttt{master} — 2026-09-07}
\end{center}

\end{document}
